% !TeX root = ../../main.tex
\subsection{疑問詞の種類}
    疑問詞はそもそも次の8つである。

    \begin{definitionbox}{\textbf{疑問詞の種類}}
        \begin{enumerate}
            \item \makebox[7em][l]{\fitblankbf[]{what}}\hbox to 3em{\dotfill\hss}\makebox[11em][l]{\fitblankbf[]{何を}}
            $\Longrightarrow$ 大きく \fitblankbf[]{もの} を問う
            \item \makebox[7em][l]{\fitblankbf[]{who}}\hbox to 3em{\dotfill\hss}\makebox[11em][l]{\fitblankbf[]{誰が}}
            $\Longrightarrow$ \fitblankbf[]{人} を問う
            \item \makebox[7em][l]{\fitblankbf[]{where}}\hbox to 3em{\dotfill\hss}\makebox[11em][l]{\fitblankbf[]{どこで}}
            $\Longrightarrow$ \fitblankbf[]{場所} を問う
            \item \makebox[7em][l]{\fitblankbf[]{when}}\hbox to 3em{\dotfill\hss}\makebox[11em][l]{\fitblankbf[]{いつ}}
            $\Longrightarrow$ \fitblankbf[]{時} を問う
            \item \makebox[7em][l]{\fitblankbf[]{whose}}\hbox to 3em{\dotfill\hss}\makebox[11em][l]{\fitblankbf[]{誰の（もの）}}
            $\Longrightarrow$ \fitblankbf[]{所有者} を問う
            \item \makebox[7em][l]{\fitblankbf[]{which}}\hbox to 3em{\dotfill\hss}\makebox[11em][l]{\fitblankbf[]{どの}}
            $\Longrightarrow$ \fitblankbf[]{選択肢} を問う
            \item \makebox[7em][l]{\fitblankbf[]{why}}\hbox to 3em{\dotfill\hss}\makebox[11em][l]{\fitblankbf[]{なぜ}}
            $\Longrightarrow$ \fitblankbf[]{理由や目的} を問う
            \item \makebox[7em][l]{\fitblankbf[]{how}}\hbox to 3em{\dotfill\hss}\makebox[11em][l]{\fitblankbf[]{どのように}}
            $\Longrightarrow$ \fitblankbf[]{方法や手段} を問う
        \end{enumerate}
    \end{definitionbox}

\subsection{特別な疑問詞の使い方}
    特別な疑問詞の使い方については、大きく分けて次のようなものがある。

    \begin{enumerate}[itemsep=0.8em]
        \item \fitblankbf[]{疑問詞 + 名詞}で1つの疑問詞として扱う使い方
        \item \fitblankbf[]{疑問詞 + 形容詞 [ 副詞 ]}で1つの疑問詞として扱う使い方
        \item 疑問詞が \fitblankbf[]{主語}として使われる使い方
        \item \fitblankbf[]{感嘆文}で使われる使い方
    \end{enumerate}

    それぞれについて、次の小節で深掘りしていくこととする。
\newpage

\subsection{(疑問詞+名詞)で1つの疑問詞として考える}
    (疑問詞+名詞) で1つの疑問詞として扱える疑問詞は\fitblankbf[]{what}、\fitblankbf[]{which}、\fitblankbf[]{whose}の3つのみである。

    \begin{theorembox}{\textbf{(疑問詞+名詞)で1つの疑問詞として考えるとき}}
        (疑問詞+名詞)で1つの疑問詞として \fitblankbf[]{文頭} に置き、その後は \fitblankbf[]{疑問文の形} を続ける
    \end{theorembox}
    \vspace{1em}

    \begin{enumerate}
        \item \fitblankbf[]{what + 名詞}のとき
        \par\vspace{0.5em}
        \noindent
        \begin{minipage}[t]{0.48\linewidth}
            \begin{itemize}
                \item 何時〜
                \dialogueblankcolumn{What time 〜}
            \end{itemize}
        \end{minipage}\hfill
        \begin{minipage}[t]{0.48\linewidth}
            \begin{itemize}
                \item どんな本〜
                \dialogueblankcolumn{What book 〜}
            \end{itemize}
        \end{minipage}
        \par\vspace{2em}
        \item \fitblankbf[]{which + 名詞}のとき
        \par\vspace{0.5em}
        \noindent
        \begin{minipage}[t]{0.48\linewidth}
            \begin{itemize}
                \item どの道〜
                \dialogueblankcolumn{Which way 〜}
            \end{itemize}
        \end{minipage}\hfill
        \begin{minipage}[t]{0.48\linewidth}
            \begin{itemize}
                \item どの本〜
                \dialogueblankcolumn{Which book 〜}
            \end{itemize}
        \end{minipage}
        \par\vspace{2em}
        \item \fitblankbf[]{whose + 名詞}のとき
        \par\vspace{0.5em}
        \noindent
        \begin{minipage}[t]{0.48\linewidth}
            \begin{itemize}
                \item 誰のかばん〜
                \dialogueblankcolumn{Whose bag 〜}
            \end{itemize}
        \end{minipage}\hfill
        \begin{minipage}[t]{0.48\linewidth}
            \begin{itemize}
                \item 誰のペン〜
                \dialogueblankcolumn{Whose pen 〜}
            \end{itemize}
        \end{minipage}
        \par\vspace{0.5em}
    \end{enumerate}
\vspace{1em}

\subsection{(疑問詞+形容詞 [ 副詞 ])で1つの疑問詞として考える}
    (疑問詞+形容詞 [ 副詞 ]) で1つの疑問詞として扱える疑問詞は\fitblankbf[]{how}のみである。

    \begin{theorembox}{\textbf{(疑問詞+形容詞 [ 副詞 ])で1つの疑問詞として考えるとき}}
        (疑問詞+形容詞 [ 副詞 ])で1つの疑問詞として \fitblankbf[]{文頭} に置き、その後は \fitblankbf[]{疑問文の形} を続ける
    \end{theorembox}
    \vspace{1em}

        \begin{minipage}[t]{0.48\linewidth}
            \begin{itemize}
                \item どのくらいの期間〜
                \dialogueblankcolumn{How long 〜}
                \item どのくらいの距離〜
                \dialogueblankcolumn{How far 〜}
                \item いくつの〜
                \dialogueblankcolumn{How many + 名詞の複数形 〜}
                \item いくら〜
                \dialogueblankcolumn{How much 〜}
            \end{itemize}
        \end{minipage}\hfill
        \begin{minipage}[t]{0.48\linewidth}
            \begin{itemize}
                \item どのくらいよく（頻度）〜
                \dialogueblankcolumn{How often 〜}
                \item 何歳〜
                \dialogueblankcolumn{How old 〜}
                \item 〜はどうですか
                \dialogueblankcolumn{How about 〜?}
            \end{itemize}
        \end{minipage}
\newpage

\subsection{疑問詞が主語として使われる}
    疑問詞が主語として使われる疑問詞は\fitblankbf[]{what}、\fitblankbf[]{who}の2つのみである。

    \begin{theorembox}{\textbf{疑問詞が主語として使われるとき}}
        疑問詞が主語として使われるときは、\fitblankbf[]{文頭} に置き、その後は \fitblankbf[]{肯定文の形} を続ける

        この際、文末の \fitblankbf[]{?} を忘れないこと。

        また、主語は \fitblankbf[]{三人称・単数} 扱いである。
    \end{theorembox}
    \vspace{1em}

    \begin{enumerate}
        \item \fitblankbf[]{what}が主語のとき
        \par\vspace{0.5em}
        \noindent
        \begin{minipage}[t]{0.48\linewidth}
            \begin{itemize}
                \item 何があなたを幸せにするの？
                \dialogueblankcolumn{What makes you happy?}
                \item[\answermark] Music does.
                \dialogueblankcolumn{音楽だよ}
            \end{itemize}
        \end{minipage}\hfill
        \begin{minipage}[t]{0.48\linewidth}
            \begin{itemize}
                \item What scared the cat?
                \dialogueblankcolumn{何がその猫を怖がらせたの？}
                \item[\answermark] 全部だよ。
                \dialogueblankcolumn{Everything did.}
            \end{itemize}
        \end{minipage}
        \par\vspace{2em}
        \item \fitblankbf[]{who}が主語のとき
        \par\vspace{0.5em}
        \noindent
        \begin{minipage}[t]{0.48\linewidth}
            \begin{itemize}
                \item 誰がここに住んでいるの？
                \dialogueblankcolumn{Who lives here?}
                \item[\answermark] ケンとジョンだよ。
                \dialogueblankcolumn{Ken and John do.}
            \end{itemize}
        \end{minipage}\hfill
        \begin{minipage}[t]{0.48\linewidth}
            \begin{itemize}
                \item Who broke the window yesterday?
                \dialogueblankcolumn{誰が昨日その窓を壊したの？}
                \item[\answermark] リクが \dots
                \dialogueblankcolumn{Riku did.}
            \end{itemize}
        \end{minipage}
        \par\vspace{0.5em}
    \end{enumerate}
\vspace{2em}

\subsection{特別な答え方をする疑問文}
    Why から始まる疑問文は、答え方が \fitblankbf[]{2パターン} ある。

    \begin{enumerate}
        \item 接続詞 \fitblankbf[]{because}を使うとき
        \par\vspace{0.5em}
        \noindent
        \begin{minipage}[t]{0.48\linewidth}
            \begin{itemize}
                \item Why are you happy?
                \dialogueblankcolumn{なんで幸せなの？}
            \end{itemize}
        \end{minipage}\hfill
        \begin{minipage}[t]{0.48\linewidth}
            \begin{itemize}
                \item[\answermark] Because I passed the exam.
                \dialogueblankcolumn{試験に合格したからだよ。}
            \end{itemize}
        \end{minipage}
        \par\vspace{2em}
        \item \fitblankbf[]{不定詞}の副詞的用法を使うとき
        \par\vspace{0.5em}
        \noindent
        \begin{minipage}[t]{0.48\linewidth}
            \begin{itemize}
                \item なんで昨日図書館に行ったの？
                \dialogueblankcolumn{Why did you go to the library yesterday?}
            \end{itemize}
        \end{minipage}\hfill
        \begin{minipage}[t]{0.48\linewidth}
            \begin{itemize}
                \item[\answermark] 本を借りるためだよ。
                \dialogueblankcolumn{To borrow some books.}
            \end{itemize}
        \end{minipage}
        \par\vspace{0.5em}
    \end{enumerate}
\vspace{2em}
